\section{Related Work}
\label{sec:related}

% Target length: ~half a page (3 short paragraphs). Organize by
% methodology, not paper-by-paper.

\paragraph{Vision Transformers and small-data training.}
Vision Transformer (ViT) showed that a pure transformer can work well for image classification when it is trained at large scale and then transferred to downstream tasks \citep{dosovitskiy2021vit}. At the same time, its original results also suggested that data and training recipe matter a lot, especially when the target task is smaller than the large-scale pre-training setup. DeiT later showed that ViTs can be trained in a more data-efficient way, but only with a strong recipe that combines several regularization and augmentation components \citep{touvron2021deit}. In this sense, convolutional networks remain a useful baseline in small-data settings, and our study revisits this issue in the fine-tuning regime on CIFAR-100.

\paragraph{Regularization in modern training pipelines.}
Modern image classification pipelines rarely rely on a single regularizer. Instead, they usually combine AdamW for decoupled weight decay \citep{loshchilov2019adamw}, dropout or stochastic depth for model regularization \citep{srivastava2014dropout,huang2016stochasticdepth}, label smoothing for softer targets \citep{szegedy2016rethinking,muller2019labelsmoothing}, and a strong augmentation suite such as RandAugment, Mixup, and CutMix \citep{cubuk2020randaugment,zhang2018mixup,yun2019cutmix}. These choices often work well together, but they are usually tuned as one package, which makes it hard to tell which part is responsible for a gain. This motivates a controlled comparison where we vary one regularization axis at a time.

\paragraph{Scope of this work.}
This paper is a controlled empirical study rather than a novel-method paper. We do not introduce a new model or a new regularizer; instead, we keep the backbone, dataset, and training budget fixed, and compare different regularization choices under the same setting. Our goal is to better understand how these methods behave for ViT fine-tuning on CIFAR-100, rather than to claim a new state-of-the-art result.


% \paragraph{Vision Transformers and small-data training.}
% % 1 paragraph: ViT \citep{dosovitskiy2021vit} matches CNNs at scale but
% % trails them when trained from scratch on small datasets. DeiT
% % \citep{touvron2021deit} closes much of this gap with a heavy
% % regularization and augmentation recipe (RandAugment, Mixup, CutMix,
% % stochastic depth, label smoothing), demonstrating that for ViTs
% % \emph{regularization is not optional}. We revisit this lesson in the
% % fine-tuning regime on CIFAR-100.
% \draftnote{3--4 sentences: ViT, DeiT, ConvNets-as-baselines context.}

% \paragraph{Regularization in modern training pipelines.}
% % 1 paragraph: Modern recipes routinely combine several regularizers:
% % decoupled weight decay in AdamW \citep{loshchilov2019adamw}, dropout
% % \citep{srivastava2014dropout} and stochastic depth
% % \citep{huang2016stochasticdepth}, label smoothing
% % \citep{szegedy2016rethinking,muller2019labelsmoothing}, and an
% % augmentation suite of RandAugment \citep{cubuk2020randaugment},
% % Mixup \citep{zhang2018mixup}, and CutMix \citep{yun2019cutmix}. These
% % are usually tuned jointly, which makes their individual contributions
% % unclear---the gap our controlled study is meant to fill.
% \draftnote{3--4 sentences covering AdamW, dropout/droppath, label
% smoothing, and the modern aug suite. Cite the keys above.}

% \paragraph{Scope of this work.}
% % 1 short paragraph: We do not propose a new regularizer or model.
% % This is a controlled empirical study under a single fixed budget,
% % varying one regularization axis at a time. Our goal is comparative
% % understanding for ViT in the small-data fine-tuning regime, not a
% % claim of state of the art.
% \draftnote{Frame the paper as a controlled empirical study, not a
% novel-method paper. This sets reviewer expectations correctly.}
